% ==============================================================================
\section{Problem Statement}
\label{sec:problem}
% ==============================================================================

We model a hangar as a directed graph $G = (\calP, \calA)$, where
$\calP = \{p_1, \dots, p_P\}$ is the finite set of \emph{maintenance positions}
and $\calA \subseteq \calP \times \calP$ is a set of \emph{blocking arcs}.  An
arc $(p, p') \in \calA$ means that position $p$ (\emph{front}) physically
blocks access to position $p'$ (\emph{rear}): an aircraft at position $p'$
cannot enter or leave the hangar without temporarily displacing the aircraft
at position $p$.  In practice $G$ is acyclic, so a position cannot block
itself transitively, and we define the \emph{blocking load} of a position as
$\ell(p) = |\{q \in \calP : p \text{ can reach } q \text{ in } G\}|$.  A
higher blocking load means that aircraft at $p$ can obstruct more other
positions, making $p$ topologically \emph{central}.

\begin{definition}[Blocking arc]
  An arc $(p, p') \in \calA$ induces an \emph{entry movement} whenever an
  aircraft arrives at position $p'$ while another aircraft is being serviced
  at position $p$, and an \emph{exit movement} whenever an aircraft departs
  from position $p'$ while another aircraft is still at position $p$.  Each
  movement contributes 2 to the movement count, one for the displacement of
  the front aircraft and one for its return.
\end{definition}

Let $\calR = \{r_1, \dots, r_R\}$ be the set of aircraft.  Each aircraft
$r \in \calR$ comes with an ordered set of maintenance jobs
$\calJ_r \subseteq \calJ$, where $\calJ$ is the union of all job sets and
each job $j \in \calJ_r$ has a fixed processing time $D_j > 0$.  The jobs of
each aircraft are linked by precedences $\calE_r \subseteq \calJ_r \times
\calJ_r$, where $(j, j') \in \calE_r$ means that job $j$ must finish before
job $j'$ starts; the pair $(\calJ_r, \calE_r)$ forms a DAG and we denote by
$j^r_1$ and $j^r_L$ its unique first and last jobs.  Each aircraft has an
\emph{earliest start} $E_r \ge 0$ before which processing cannot begin and a
\emph{target finish} $L_r > E_r$ beyond which completion incurs delay.  The
\emph{total processing time} of aircraft $r$ is the parameter
$D_r = \sum_{j \in \calJ_r} D_j$.

A solution to the AHPP makes two coupled decisions.  At the assignment
level, it sends each aircraft $r$ to exactly one position $\pi(r) \in \calP$.
At the scheduling level, it fixes a start time $t^S_j \ge 0$ for every job
$j \in \calJ$, from which the aircraft start time $s_r = t^S_{j^r_1}$ and
the aircraft finish time $f_r = t^S_{j^r_L} + D_{j^r_L}$ follow.

Any solution must satisfy the following requirements.  Every aircraft is
\emph{assigned to exactly one position}, and a position handles its
aircraft \emph{sequentially}: when two aircraft $r$ and $r'$ are assigned
to the same position, their occupancy intervals
$[s_r, f_r]$ and $[s_{r'}, f_{r'}]$ must not overlap and a minimum
separation $\varepsilon > 0$ must elapse between the finish of one and the
start of the next.  Within each aircraft,
jobs must respect their \emph{precedence chain}, so $t^S_{j'} \ge t^S_j +
D_j$ for every $(j, j') \in \calE_r$, and the aircraft's start time must
respect its \emph{earliest-start time window}, $s_r \ge E_r$.

The blocking-arc structure imposes the final family of requirements.  For
every arc $(p, p') \in \calA$ and every pair of aircraft $r, r'$ with
$\pi(r) = p$ and $\pi(r') = p'$, an entry movement is incurred if $r'$
arrives at the rear position while $r$ is still being serviced at the
front, that is, if $s_r \le s_{r'} \le f_r - \varepsilon$; analogously,
an exit movement is incurred whenever $r'$ departs from the rear while
$r$ is still at the front, that is, if $s_r \le f_{r'} \le f_r -
\varepsilon$.  Each such event triggers a temporary relocation of the
front aircraft and is counted as two movements (displacement plus return).

\paragraph{Single $\varepsilon$ for two roles.}
The parameter $\varepsilon$ plays the same physical role in both of the
requirements above: it is the time needed to tow an aircraft in or out of
a hangar slot.  This primitive operation has the same duration whether it
is performed as a turnaround between two consecutive maintenance slots at
one position --- the role of $\varepsilon$ in the separation requirement
--- or as a temporary displacement of the front aircraft to free a
blocked rear position --- the role of $\varepsilon$ in the
$f_r - \varepsilon$ margin of the blocking conditions.  We therefore
treat the two appearances as a single parameter; for the reported
benchmark we use $\varepsilon = 0.5$ days, an indicative half-day
allowance for this primitive towing operation, while all other temporal
magnitudes are expressed in integer days.  The value is recorded in each
instance file, so that experiments on hangars with different operational
rhythms can use a different $\varepsilon$ without changing the model.

The objective is a weighted combination of three criteria,
\begin{equation}
  \label{eq:obj}
  \min \;\; W^M \cdot m
        \;+\; W^D \cdot \sum_{r \in \calR} v^D_r
        \;+\; W^S \cdot n,
\end{equation}
where $m = \max_{j \in \calJ}(t^S_j + D_j)$ is the makespan,
$v^D_r = \max(0,\, f_r - L_r)$ is the delay of aircraft $r$, $n$ counts
the blocking movements induced above, and $W^M, W^D, W^S \ge 0$ are
user-defined weights.

The AHPP is $\mathcal{NP}$-hard, since it subsumes the single-machine
weighted-completion-time scheduling problem with arbitrary release dates,
which is already $\mathcal{NP}$-hard~\citep{Lenstra1977}; the additional
position-assignment layer and the blocking-arc constraints further enlarge
the search space.  For a benchmark instance with $P = 5$ positions and
$R = 10$ aircraft (each with $4$--$6$ jobs), the MILP formulation of
\S\ref{sec:milp} involves roughly $1{,}000$ binary variables and $3{,}000$
constraints; at $R = 30$ these grow by an order of magnitude, motivating
the heuristic approaches of \S\ref{sec:approaches}.
